% Auto-generated by scripts/make_table1.py -- do not hand-edit.
% Requires: \usepackage{booktabs} \usepackage{graphicx}
% A five-column table with prose in three of them: it needs the full
% text width of a two-column page, so it is a table* float.
\begin{table*}[t]
\centering
\caption{Nearest-neighbour work on evaluation protocols, with the part of this paper that addresses what each left open.}
\label{tab:nearest-neighbours}
\resizebox{\textwidth}{!}{%
\begin{tabular}{@{}llp{0.20\textwidth}p{0.20\textwidth}p{0.24\textwidth}@{}}
\toprule
Work & Venue & What it established & What it left open & What this paper adds \\
\midrule
\texttt{keles2021psnr} & PCS 2021 & PSNR over a set is not the mean of per-image PSNR; gives the bound for SISR (2.07 dB measured, 2.507 dB theoretical) & Isolates aggregation only; no protocol matrix, and no seed variance & Protocol P1 (best-of-N) is the extreme case of the same effect; we quantify it against seven other protocol dimensions \\
\texttt{parmar2022cleanfid} & CVPR 2022 & Antialiased vs aliased resizing changes FID rankings in GAN evaluation & Resize semantics only, in the perceptual/GAN metric setting & P6 vs P5 isolates antialiasing in the distortion metric; the effect is 2.72 dB, not a ranking detail \\
\texttt{zhang2024seal} & ICLR 2024 & Random-crop evaluation is unreliable and flips rankings among SR methods & Sampling randomness only; degradation, channel and aggregation are held fixed & P1 vs P2 vs P3 separates crop position from the other five dimensions, with 5 seeds per cell \\
\texttt{srinivasan2026lol} & arXiv 2026 & LOL benchmarks carry hidden label shifts: LOLv1 test labels render differently, and LOLv2-Real is 14\% darker than the training labels & Dataset-side label shift; the evaluation protocol itself is not varied & The low-light task is our counterexample: protocol does not dominate there (ratio 0.21), so the claim is conditional \\
\texttt{shah2026hyperbench} & arXiv 2026 & Protocol implementation differences in the Wald protocol prevent hyper-resolution rankings from holding & One task family (hyperspectral), one protocol dimension (degradation) & Three tasks and eight protocols; the protocol span exceeds the training gain on two of them \\
\texttt{su2025bridging} & arXiv 2025 & Audits perceptual quality metrics for SR and finds they diverge from distortion metrics & Metric choice is orthogonal to the protocol dimensions this paper varies & Holds the metric fixed and varies sampling, channel, aggregation and degradation instead \\
\texttt{cawley2010overfitting} & JMLR 2010 & Selection bias: repeated model selection on a validation set inflates the reported performance of the selected model & Gives the theory; does not measure how large the bias is under a stated protocol & Measures the same bias at 1.74 dB under P1, which is why the reporting checklist requires the selection budget N \\
\bottomrule
\end{tabular}}
\end{table*}
